% =========================
% Chapter 5
% =========================
\chapter{Task Families and Algorithms}
\label{ch:task_families}

This chapter details the three task families. Throughout, a common structure recurs:
a perception stage that turns frames into \emph{trusted} detections, a pure decision
layer that turns detections into guidance, and a speech policy that decides what
actually reaches the user's ears. The perception stages differ per task; the
philosophy --- nothing unverified is ever spoken --- does not.

\section{Object Allocation}
\label{sec:object_allocation}

Object Allocation answers ``find my cup''. Its loop runs at 5\,FPS over the
session's frame stream.

\subsection{Detection and Temporal Filtering}

Each frame passes through YOLOv8n restricted to the requested class, with a 0.35
confidence floor. Raw per-frame detections are too flickery to speak: a
\emph{temporal-consistency filter} requires the target to appear in at least 3 of
the last 5 frames before it is trusted, killing single-frame false positives. When
multiple instances are visible (two chairs), guidance targets the most head-on one
--- the box whose centre is closest to the frame centre.

\subsection{Spatial Reasoning with Per-Class Size Priors}
\label{sec:spatial_reasoning}

The trusted box is classified into a horizontal region --- left / center / right,
split at 0.35 and 0.65 of frame width --- and a distance bucket. Distance uses
the larger of the box's width and height fractions, compared against
\emph{per-class thresholds}: a laptop at 40\% of frame width is near, a cup at 18\%
is equally near, a phone at 12\% likewise --- same arithmetic, three correct
answers, because the priors encode real object sizes. Bucket transitions are
debounced (two consecutive frames of the same bucket) so YOLO's box jitter at a
threshold cannot fire three contradictory phrases in a second; when the client's
motion monitor reports that the user is actually moving, the debounce is bypassed so
guidance keeps up with real position changes.

\subsection{Guidance Policy}

The tracker speaks on first sighting (``Found your cup, to your left''), on bucket
changes (``Your cup is straight ahead, a few steps away''), and re-affirms unchanged
guidance at most every 6\,s. Edge cases have explicit behaviors: never seen within
30\,s produces periodic scan prompts (``I don't see your cup yet, slowly turn
around''); a target that vanishes for a full temporal window produces a one-shot
loss announcement including its last known direction; never seen within 60\,s aborts
the task with an apology. Completion is user-driven (``got it'') --- with one
exception, below.

\section{Reach Guidance}
\label{sec:reach_guidance}

When the target's distance bucket is \emph{near} and MediaPipe detects a hand in
frame, guidance switches from body-relative to \emph{hand-relative} cues. MediaPipe
runs only in this phase --- invoking it earlier would waste computation while the
target is across the room.

The controlling quantity is the offset between the index fingertip and the target
box centroid, normalized by frame size. Three states result: \emph{approach}
(directional cues: ``Move your hand to the right'', ``Raise your hand up''),
\emph{almost} (fingertip within 10\% of the target centroid: ``Almost there. Reach
forward.''), accompanied by a short haptic pulse, and \emph{touching} (fingertip
inside the target box): ``Your hand is on the cup. Grasp it,'' with a long haptic
pulse. Touching is the single autonomous success exit in the whole system --- it
fires \texttt{task\_complete} without user confirmation, because the fingertip on
the object \emph{is} the confirmation. A grace period keeps the system in reach mode
across momentary hand-detection dropouts, preventing oscillation between guidance
modes at 5\,FPS.

\section{Navigation: Goal-Directed Exploration}
\label{sec:navigation}

Navigation is the deepest task family. The user names only a destination room; the
system explores toward it with no map, no markers, and no localization --- the
design bet is that \emph{semantic room recognition} plus \emph{door detection} plus
\emph{occasional human hints} can substitute for the SLAM a phone camera cannot do.

\subsection{Destination Model and Semantic Arrival}
\label{sec:goals_model}

Each supported destination maps to \emph{indicator objects}, split into primary
(near-unique to the room) and secondary (supporting, often shared across rooms), as
shown in Table~\ref{tab:goal_indicators}. Arrival is declared when at least one
primary or at least two secondary indicators are confirmed --- so a lone sink never
declares a kitchen (bathrooms have sinks), but a refrigerator does. All indicators
are COCO classes, so the stock detector produces them with no additional training.

\begin{table}[H]
\centering
\caption[Destination rooms and indicator objects]{Destination rooms and their indicator objects (all COCO classes). Arrival
requires ${\geq}1$ primary or ${\geq}2$ secondary indicators, temporally confirmed.}
\label{tab:goal_indicators}
\begin{tabular}{lll}
\toprule
\textbf{Goal} & \textbf{Primary indicators} & \textbf{Secondary indicators} \\
\midrule
kitchen      & refrigerator, oven, microwave, toaster & sink, dining table \\
bathroom     & toilet                                  & sink \\
bedroom      & bed                                     & --- \\
living room  & couch, TV                               & potted plant, remote \\
office       & laptop, keyboard                        & mouse, TV \\
dining room  & dining table                            & chair \\
\bottomrule
\end{tabular}
\end{table}

Indicator evidence is deliberately \emph{localized}: during a room scan, a class is
confirmed only when enough sightings fall within one $90^\circ$ window. A real
fridge racks up its sightings in one spot as the camera sweeps past; detector noise
scatters around the room and can never sum to an arrival.

\subsection{The Door Detector}
\label{sec:door_model}

COCO has no door class, so a dedicated detector was trained. Two public datasets
were merged and reduced to a single \texttt{door} class: DoorDetect
\cite{arduengo2021door} (386 door images plus 827 hard negatives --- handles and
cabinets --- that teach the model what a door is \emph{not}) and DeepDoors2
\cite{ramoa2021deepdoors} (3{,}000 images with 3{,}106 door boxes, converted from
segmentation masks), for a combined 4{,}213 images split 80/20 into 3{,}371 training
and 842 validation images. A YOLOv8n was fine-tuned from COCO-pretrained weights
(319 of 355 layers transferred) at $640$ input resolution, batch 8, AdamW; the best
checkpoint (epoch 39) achieves \textbf{mAP@50 $= 0.946$}, mAP@50--95 $= 0.829$,
precision $0.968$, and recall $0.892$ --- comfortably above the ${\geq}0.6$ mAP@50
target set for the milestone. A second, 4-class model (door / handle / cabinet door
/ refrigerator door) was trained on the original DoorDetect labels; its recall is
too low to detect doors alone, but that is not its job --- it serves as a
\emph{semantic verifier}, below.

\subsection{The Three-Layer Door Perception Funnel}
\label{sec:door_funnel}

A false door is the most dangerous percept in the system --- it points a blind user
at a wall, a curtain, or a refrigerator. Raw detector confidence cannot prevent
this: in live testing, blank walls scored as doors at 0.9 confidence. Every door
candidate therefore descends a three-layer funnel (Figure~\ref{fig:door_funnel}),
each layer motivated by a specific observed failure.

\begin{figure}[H]
\centering
\resizebox{0.98\textwidth}{!}{%
\begin{tikzpicture}[node distance=5mm]
  \node[box, minimum width=118mm] (cand) {Door model candidates (conf $\geq 0.45$)};
  \node[box, minimum width=118mm, below=of cand] (geo)
    {\textbf{Layer 1 --- geometric gates:} aspect ratio $\leq 1.4$ (wall spans);
     frame-fill $\leq 0.80$ during scan (whole-scene latches);\\
     Canny edge density (blank walls) --- smooth candidates \emph{demoted}, not dropped};
  \node[box, minimum width=118mm, below=of geo] (sem)
    {\textbf{Layer 2 --- semantic verification (4-class model):} strong textured
     candidates pass alone;\\ weak or smooth candidates require corroboration ---
     verifier sees a \emph{door} or a \emph{handle} there\\ (kills curtains);
     verifier's \emph{refrigerator door} without door evidence vetoes};
  \node[box, minimum width=118mm, below=of sem] (arb)
    {\textbf{Layer 3 --- fridge arbitration:} a COCO ``refrigerator'' overlapping a
     door candidate\\ is kept only if the verifier confirms a refrigerator door ---
     otherwise it \emph{is} the door\\ (prevents false kitchen arrivals from
     door-as-fridge confusions)};
  \node[box, minimum width=118mm, below=of arb] (rel)
    {\textbf{Reliability gates:} $\geq 3$ sighting units before a door is spoken or
     targeted\\ (verifier-corroborated frames count double); 2-of-3 temporal
     confirmation per frame};
  \draw[arr] (cand) -- (geo);
  \draw[arr] (geo) -- (sem);
  \draw[arr] (sem) -- (arb);
  \draw[arr] (arb) -- (rel);
\end{tikzpicture}}
\caption[Door perception funnel]{The door perception funnel. Each layer was added in response to a specific
live failure: whole-wall detections, dim plain doors indistinguishable from walls,
a lace curtain that beat every geometric gate, and doors misread as refrigerators.}
\label{fig:door_funnel}
\end{figure}

One subtlety deserves emphasis: edge density is \emph{not} a hard reject. A plain
door in dim light is as smooth as a wall (both ${\sim}0.01$ edge fraction), and
confidence cannot break the tie (walls reach 0.9). Smooth candidates are instead
demoted to require semantic proof --- only the verifier seeing a door or a handle
there can pass them. This is what finally separated the lace curtain (edge-rich,
tall, door-sized, defeats all geometry) from real doors: curtains have no handles.

\subsection{The Compass-Anchored Room Scan}
\label{sec:room_scan}

On entering a room (and at journey start), the system requests one slow full turn:
``Slowly turn to your right, all the way around, until you are facing where you
started.'' The first compass heading is locked as the anchor --- ``ahead'' means
where the scan began --- and every subsequent detection is filed by true bearing:
$\text{bearing} = \text{camera heading} + (x_{\text{frame}} - 0.5)\cdot
\text{FOV}$, the in-frame correction fixing a class of ``said left, was right''
errors. Sightings are accumulated into $12 \times 30^\circ$ sectors, and door
bearings are \emph{clustered} (sightings within $30^\circ$ merge, including across
the $\pm180^\circ$ seam) so one physical door is announced once, with one precise
heading.

Scan completion requires \emph{both} ${\sim}350^\circ$ of accumulated rotation
\emph{and} a compass-confirmed return to within $25^\circ$ of the start heading. An
earlier design that ended on sector coverage alone could be faked by compass noise
mid-turn, ending the scan while the user faced the wrong way --- and every spoken
direction would then be computed from a broken premise. With the current rule, the
directions in the scan summary (``I found a door behind you, to the right'') are
true at the moment they are heard. A motion gate skips detection on blurred frames
while the phone pans too fast (coaching ``slow down'' only after several consecutive
fast frames), but keeps advancing the rotation total so a fast sweep cannot stall
completion. Without a compass (e.g., a laptop), the scan falls back to a single
timed steady pass with no bearing summaries.

\subsection{The Exploration Controller}
\label{sec:exploration_controller}

\begin{figure}[H]
\centering
\resizebox{0.95\textwidth}{!}{%
\begin{tikzpicture}[node distance=8mm and 10mm]
  \node[state, minimum width=34mm] (disc) {\textbf{discover}\\silent guided $360^\circ$ scan};
  \node[state, minimum width=30mm, below left=12mm and 2mm of disc] (face) {\textbf{face\_target}\\turn user to face\\door or sighting};
  \node[state, minimum width=27mm, below right=12mm and 2mm of disc] (arr) {\textbf{arrived}\\declare, task ends};
  \node[state, minimum width=27mm, below=10mm of face] (godoor) {\textbf{go\_door}\\guide in: steps + hand cue\\obstacle watchdog ON};
  \node[state, minimum width=27mm, right=14mm of godoor] (goind) {\textbf{go\_indicator}\\directed re-confirmation};

  \draw[arr] (disc) -- node[lbl, right]{strong indicators\\(after directed check)} (arr.north);
  \draw[arr] (disc) -- node[lbl, left]{door chosen /\\weak indicator} (face);
  \draw[arr] (face) -- node[lbl, left]{facing door} (godoor);
  \draw[arr] (face) -- node[lbl, above right=0 and -6mm]{facing sighting} (goind);
  \draw[arr] (goind) -- node[lbl, right]{confirmed} (arr);
  \draw[arr] (goind) to[bend right=18] node[lbl, above]{not confirmed:\\take door / rescan} (face);
  \draw[arr] (godoor.west) to[bend left=40] node[lbl, left]{doorway transit\\$\rightarrow$ new room} (disc.west);
  \draw[arr] (godoor) to[bend right=25] node[lbl, below]{door lost 12\,s\\$\rightarrow$ rescan} (disc.south east);
\end{tikzpicture}}
\caption[Navigation controller state machine]{The navigation controller's internal state machine (nested inside the
session FSM's NavigationActive state). Priority is absolute: confirmed indicators
always beat doors --- the system never walks the user out of the room they were
looking for.}
\label{fig:nav_controller}
\end{figure}

The controller (Figure~\ref{fig:nav_controller}) implements a two-phase principle:
Phase~1 scans silently and collects evidence; Phase~2 acts on that evidence by
priority. Strong indicator evidence still receives a \emph{directed double-check} ---
the user is turned toward the sightings and the system re-confirms before declaring
arrival, so a scan artifact alone never ends a journey. Weak indicator sightings get
the same directed confirmation, and fall back to the door branch if unconfirmed.
When only doors were found, the most-sighted door (tie-broken toward straight ahead)
becomes the target; the user is rotated to face it using live compass differences,
then guided in. If nothing useful was found, the room is rescanned --- and all such
announcements are merged into single utterances so instructions never collide.

\subsection{Door Approach: Distance in Steps}
\label{sec:door_distance}

Door distance uses a pinhole model over the door's known height: with doors
${\sim}2$\,m tall and the focal length derived from the camera's field of view,
$d = h_{\text{real}} \cdot f_{\text{px}} / h_{\text{px}}$. Two corrections came from
live testing: the focal length must be derived from the image's \emph{long} side
(phones stream portrait; using width understated every distance by ${\sim}33\%$),
and when the verifier also boxed the door, distance is measured on the tighter of
the two boxes (the primary model's boxes run loose at range, and a fat box reads as
near). A one-shot calibration constant scales the whole chain from a single
ground-truth measurement. Distances are spoken in steps (${\sim}0.75$\,m each) with
a tactile endgame: within three steps, ``reach out with your hand to find the
door.'' The spoken sequence is strictly ordered --- one call-out stating where the
door is, then a path check, then either ``the path is clear, walk\ldots'' or an
obstacle warning --- so safety information always precedes the instruction to move.

\subsection{The Obstacle Watchdog}
\label{sec:obstacles}

Armed only in the walking phase (\texttt{go\_door}) --- the only phase the user
moves --- the watchdog fuses two layers with absolute speech priority over door
guidance. \textbf{Layer A (named objects):} during the approach, the COCO detector
additionally requests obstacle classes (person, chair, table, bags, pets, \ldots);
any box that sits low in frame, is large enough to matter, and overlaps the central
walking lane triggers ``There's a chair in your path. Step to your left, where it's
clear.'' \textbf{Layer B (unnamed clutter):} SegFormer floor segmentation labels
every pixel; the lane just ahead of the user's feet must remain mostly
walkable-floor. Something \emph{beside} the path leaves the floor corridor
continuous (silent); something \emph{in} the path punches a hole in it (warn, with
the sidestep direction chosen toward the flank with more visible floor). This
floor-based formulation replaced an earlier relative-depth tripwire
\cite{yang2024depthanything} that could not reliably separate beside-path from
in-path objects in cluttered rooms; the depth variant is retained as a configurable
alternative. Both layers are debounced (three consecutive frames to warn, two clear
frames to release), re-warn every ${\sim}3$\,s while blocked, and close with ``Okay,
the way ahead is clear'' --- followed, in the same utterance, by re-orientation
toward the door.

\subsection{Doorway Transit Detection}
\label{sec:transit}

Detecting ``the user just walked through'' has no single reliable signal, so the
system fuses several, each gated by the others. A tracked approach must first get
\emph{near} the door (metric distance below 2\,m, or the confirmed door having grown
to fill most of the frame height); then a saturated door-shaped box at arm's length
--- which the edge gate deliberately rejects as a door, since a panel up close is
smooth --- arms an ``at the door'' latch, but only after consecutive qualifying
frames (a single loose box mid-sidestep once spoke ``you're at the door'' from two
meters away). The user hears the full sequence in one utterance: ``You're right at
the door. Reach out with your hand, open it, walk through the doorway, and take two
or three steps into the room.'' Transit is then inferred when the door disappears
from view for ${\sim}2$\,s of observed frames \emph{and} the camera reports
sustained motion --- the camera is the odometer, and detection flicker without
movement can never fake a walk-through. A hold cap forces the inference if a new
room's own door candidates keep the latch alive after the user has clearly moved on.
On transit, the room counter increments and a fresh scan begins, anchored to the
user's new heading.

\subsection{Bounded Autonomy: Room Cap and Journey Timeout}
\label{sec:room_cap}

Without localization, an exploration in the wrong wing of a building could otherwise
continue indefinitely. Two mechanisms bound it. After a configurable number of rooms
(four, then every second room), the transit announcement becomes a check-in:
``You're through --- that's 4 rooms now and no kitchen yet. Say stop anytime if
you'd like to give up, or keep going\ldots'' --- informing rather than interrogating,
since ``stop'' is already honored globally. Independently, a whole-journey timeout
(five minutes) ends the task with a spoken explanation and an invitation to retry.
Both keep the user in control of an exploration whose length perception alone cannot
predict.

\subsection{Speech Discipline}
\label{sec:speech_discipline}

A blind user cannot glance at a transcript; overlapping or stale audio is a
functional bug. Three rules are enforced throughout navigation: \emph{priority}
lines (scan summaries, phase changes, obstacle warnings, arrival) are always spoken;
\emph{ambient} nudges (keep scanning, distance countdowns) are dropped --- never
queued --- when speech is already playing, when they repeat too soon, or hot on the
heels of another line; and announcements that would otherwise be adjacent are merged
into single utterances, because two back-to-back priority lines would cut each other
off. The scan-summary sentence exemplifies all three: completion confirmation,
findings by direction, and the next instruction arrive as one connected utterance.
